\section{Central rigidity from relative point fields}\label{sec:central-ramification}
This section proves Theorems~\ref{thm:central-p2-rigidity} and \ref{thm:ss-prime-range}. The first argument compares a cyclic extension of point fields with the residual ramification of the quotients by a canonical kernel and its transported image. The second combines that criterion with the pointwise graph valuation bound, losing at most one torsion level.

Pass to the completed maximal unramified base and write $e=v_K(p)$. Let $L/K$ be the common tame good-reduction field of degree $m=p^2-1$, with algebraically closed residue field $k$ and $E=v_L(p)=me$. Lemma~\ref{lem:tame-good} supplies this field. For a two-slope model write $h=(p-1)z$, so
\begin{equation}\label{cr-eq:AB}
 A=(p+1)e-z,\qquad B=z/p,\qquad 0<z<pe.
\end{equation}
The deformation weight makes $z$ integral. Lemmas~\ref{rc-lem:newton} and \ref{fc-lem:rational} give the division valuations and the rationality of the canonical points. The canonical line and its quotient character are trivial over $L$, since the determinant over the good field is cyclotomic and $\mu_p\subset L$.

We use the following exact quotient data, proved in Proposition~\ref{rc-prop:quotients} and the marking calculation preceding Proposition~\ref{fc-prop:central}. A canonical quotient has normalized Hasse value $pz$ for $z<e$, value $(p+1)e-z$ for $z>e$, and one slope for $z=e$. Its dual kernel is noncanonical in the first case and canonical in the second. A quotient by a noncanonical stable line has normalized Hasse value $z/p$. A one-slope model has all nonzero $p$-point values equal to $e$, all primitive $p^2$-point values equal to $e/p^2$, and every $p$-quotient has normalized Hasse value $e$.

If a $p^2$-torsion isomorphism carries $C$ to $C'$, it induces
\begin{equation}\label{cr-eq:quotient}
 (G/C)[p]\simeq(H/C')[p],\qquad
 (G/C)[p]=\{x\in G[p^2]:px\in C\}/C.
\end{equation}
The quotients exist over $K$ for $K$-stable kernels and preserve special inertia traces; see Lemma~\ref{lem:actual-quotients}. These facts require no extension of the given torsion isomorphism to good integral models.

\subsection{The cyclic relative extension and its break}
The elementary valuation estimate below fixes the lower-numbering convention; compare \citet[Chapter~IV, \S1]{Serre1979}.

\begin{lemma}\label{cr-lem:valuation}
Let $F/F_0$ be a totally ramified cyclic extension of degree $p$ of
complete discretely valued fields of characteristic zero, with common
perfect residue field of characteristic $p$. Let $w$ be the integer
valuation of $F$, let $\sigma$ be nonidentity, and let $b$ be its lower
ramification break. For every $x$ in the maximal ideal of $\mathcal O_F$,
\[
 w(\sigma(x)-x)\geq w(x)+b.
\]
Equality holds when $p\nmid w(x)$.
\end{lemma}

\begin{proof}
Choose a uniformizer $\varpi$ of $F$. The residue of
$\sigma(\varpi)/\varpi$ has order dividing $p$ in the multiplicative
group of a field of characteristic $p$, so it is~$1$. Thus $b\geq1$, and
\[
 \sigma(\varpi)=\varpi+u\varpi^{b+1}
                  +\text{terms of larger valuation},\qquad u\in\mathcal O_F^*.
\]
Use Teichm\"uller representatives from the coefficient ring of $F_0$ for
the common residue field; they are fixed by $\sigma$. Every element of
$\mathcal O_F$ has a convergent expansion in powers of $\varpi$ with
these coefficients. The binomial expansion gives
\[
 w\bigl(\sigma(\varpi^j)-\varpi^j\bigr)\geq j+b\qquad(j\geq1).
\]
When $p\nmid j$, the first binomial term has exactly this valuation, and
all later terms have larger valuation because $b\geq1$. Apply these
estimates to the expansion of $x$. If its first exponent $w(x)$ is prime
to $p$, its change has strictly smaller valuation than the changes of all
subsequent terms. This proves the equality assertion as well.
\end{proof}

\begin{proof}[Proof of Theorem~\ref{thm:central-p2-rigidity}]
Let $C$ be the canonical $p$-subgroup of $G$, choose
$0\ne P\in C(L)$, and choose $Q$ with $[p]Q=P$. The point $Q$ has exact
order $p^2$. If $T$ is the good formal coordinate, centrality gives
\begin{equation}\label{cr-eq:qval}
 v_L(T(Q))=\frac{A}{p^2}.
\end{equation}
The denominator is $p^2$, because $p\nmid A$. Thus the ramification
degree of $L(Q)/L$ is at least $p^2$. Its field degree is at most $p^2$,
since the orbit of $Q$ lies in the set of $p^2$ preimages of the fixed
point $P$. Therefore
\begin{equation}\label{cr-eq:field1}
 F=L(Q),\qquad [F:L]=p^2.
\end{equation}
The extension is totally ramified, since $k$ is algebraically closed.

Form $q:G\to G_1=G/C$. Its normalized coordinate is, up to a unit, the
product of $T-_G R$ over $R\in C$. At $Q$, all $p$ factors have valuation
$A/p^2$, since the nonzero kernel points have the larger valuation $A$.
Thus
\[
 v_L\bigl(T_1(q(Q))\bigr)=\frac Ap.
\]
The point $q(Q)$ is nonzero $p$-torsion. The normalized Hasse value of
$G_1$ is $A$, and this is a noncanonical point of $G_1[p]$. Its valuation
forces field degree at least $p$. Its orbit has at most $p$ elements,
because the canonical line of $G_1[p]$ and its quotient character are
trivial over $L$. Hence
\begin{equation}\label{cr-eq:field0}
 F_0=L(q(Q)),\qquad [F_0:L]=p.
\end{equation}

The extension $F/F_0$ has degree $p$ and is cyclic Galois. Indeed the
$p$ points $Q+R$, $R\in C(L)$, are the full fiber of $q$ above $q(Q)$ and
all belong to $F$. The formal addition series converges in the complete
field $F$. The minimal polynomial of $T(Q)$ over $F_0$ has degree $p$,
and all its roots, precisely the coordinates of this fiber, lie in $F$.
Thus a nonidentity $\sigma\in\Gal(F/F_0)$ sends $Q$ to $Q+P_0$ for some
$0\ne P_0\in C(L)$.

Normalize $w=p^2v_L$ on $F$. A formal coordinate difference under
translation by $P_0$ is $T(P_0)$ times a unit. Consequently
\[
 w(T(Q))=A,\qquad
 w\bigl(\sigma(T(Q))-T(Q)\bigr)=p^2A.
\]
Lemma~\ref{cr-lem:valuation}, with its equality clause, determines the
relative lower break:
\begin{equation}\label{cr-eq:break}
 b=(p^2-1)A.
\end{equation}

Now let $f:G[p^2]\to H[p^2]$ be the given isomorphism. Put
$C'=f(C)$, $R_Q=f(Q)$, and $q':H\to H_1=H/C'$. The line $C'$ is
$K$-stable and all its points are $L$-rational. The point stabilizers of
$Q$ and $R_Q$ in the absolute Galois group of $L$ are identical. The
isomorphism~\eqref{cr-eq:quotient} also identifies the stabilizers of
$q(Q)$ and $q'(R_Q)$. Thus, as subfields of the fixed separable closure,
\begin{equation}\label{cr-eq:transport}
 L(R_Q)=F,\qquad L(q'(R_Q))=F_0.
\end{equation}
For the same automorphism $\sigma$,
\begin{equation}\label{cr-eq:translation}
 \sigma(R_Q)=R_Q+f(P_0).
\end{equation}
These assertions use only the finite torsion isomorphism; they do not
require its extension to integral good models.

The group $G_1[p]$ has nontrivial wild action over $L$. Therefore $H$ cannot
be two-slope with $C'$ canonical in the outer range: its quotient would
have a stable noncanonical dual line and hence trivial residual action.
The canonical boundary $z_H=e$ is excluded for the same reason, since its
quotient is one-slope.

In every remaining case, $H_1$ is two-slope. Denote its normalized Hasse
value by $t$. The exact values are as follows:
\begin{enumerate}
\item If $C'$ is canonical and $H$ is central, then
      $t=(p+1)e-z_H$. The point $[p]R_Q$ is canonical, central division
      gives $v_L(S(R_Q))=t/p^2$, and $f(P_0)$ has valuation $t$.
\item If $C'$ is noncanonical, then $t=z_H/p$.
      The point $[p]R_Q$ is noncanonical, so again
      $v_L(S(R_Q))=t/p^2$, while $f(P_0)$ has valuation $t$.
      In particular $t$ is integral, since the nonzero points of $C'$
      are $L$-rational.
\item If $H$ is one-slope, then $t=e$. All nonzero $p$-points have
      valuation $e$ and every primitive $p^2$-point has valuation
      $e/p^2$, including the equality case of the Newton boundary.
\end{enumerate}
In all three cases, \eqref{cr-eq:translation} therefore gives
\[
 w(S(R_Q))=t,\qquad
 w\bigl(\sigma(S(R_Q))-S(R_Q)\bigr)=p^2t.
\]
Applying Lemma~\ref{cr-lem:valuation} to $S(R_Q)$ and using
\eqref{cr-eq:break}, we obtain
\begin{equation}\label{cr-eq:opposite}
 (p^2-1)A\leq(p^2-1)t,\qquad A\leq t.
\end{equation}

The quotient $p$-torsion representations are isomorphic, and the
normalized Hasse value $A$ is prime to $p$. If $t$ is also prime to $p$,
Theorem~\ref{rc-thm:primehasse} gives $t=A$ and equal special
trace characters of $G_1,H_1$. Their $K$-isogenous sources then have equal
special traces. If $p\mid t$, Corollary~\ref{rc-cor:strictordering} instead
gives $t<A$, contradicting~\eqref{cr-eq:opposite}. This proves the
theorem.
\end{proof}

\begin{remark}
The two ramification comparisons retain different information. The
residual representation of the quotients makes a divisible quotient
Hasse value strictly smaller than the prime-to-$p$ one. The cyclic
extension of point fields above those quotients makes it at least as
large. The latter inequality is obtained from the original
$p^2$-torsion, even though the quotient comparison itself has only
level~$p$.
\end{remark}

\subsection{The supersingular range and the full-class upper bound}
We first make explicit the pointwise estimate used in the proof. Proposition~\ref{fc-thm:valuation} gives, for an unequal-character comparison at level $p^n$,
\begin{equation}\label{cr-eq:graph}
 v_L(T(P))\ge\frac{p-1}{p^n-1}
 \quad\text{for every nonzero $P\in G[p^n]$, and likewise for $H$.}
\end{equation}
Indeed its proof obtains $x^{p^n},y^{p^n}\in\pi^{p-1}I$ in the graph algebra. Evaluating at $(P,f(P))$ and taking the smaller coordinate value $r$ gives $p^nr\ge p-1+r$, proving \eqref{cr-eq:graph}.

For a two-slope model, the minimum nonzero $p^n$-torsion value is $z/p^{2n-1}$. At level one it is $B=z/p$. The inequality $v_L([p]x)>v_L(x)$ for a nonzero formal point with nonzero $p$-multiple shows that every nonzero torsion target has valuation at most $A<E$, so the precise range of Lemma~\ref{rc-lem:newton} applies throughout the inverse chain. At every further level the minimum targets have all $p^2$ divisions on the low branch, which divides that value by $p^2$. A target of larger value below $p^2B$ has divisions of strictly larger value; if its value is at least $p^2B$, all its divisions have value at least $B$. Division points of zero are first-level points and cannot change the minimum. Induction proves the assertion. It follows that
\begin{equation}\label{cr-eq:zbound}
 z\ge\frac{(p-1)p^{2n-1}}{p^n-1}>(p-1)p^{n-1}.
\end{equation}

\begin{proof}[Proof of Theorem~\ref{thm:ss-prime-range}]
Suppose the two special traces differ. Lemma~\ref{lem:one-slope-recovery}, which has no upper bound on $e$, excludes every case involving a one-slope model. Thus both models have two slopes. If both normalized Hasse values are prime to $p$, Theorem~\ref{rc-thm:primehasse} gives equal traces already at level $p$. We may therefore choose one model $G$ whose normalized Hasse value $z$ is divisible by $p$.

If $z>e$, then $G$ is central and $A=(p+1)e-z$ is prime to $p$. Apply Theorem~\ref{thm:central-p2-rigidity} to the restricted $p^2$-isomorphism, a contradiction. Equality $z=e$ is impossible because $p\nmid e$.

Suppose $z<e$. Quotient $G$ by its canonical line $C$, and quotient $H$ by the transported $K$-stable line $f(C)$. These paired quotients have isomorphic $p^{n-1}$-torsion by Lemma~\ref{lem:actual-quotients}. The first quotient has normalized Hasse value $z_1=pz$. Equations~\eqref{eq:ss-prime-range} and \eqref{cr-eq:zbound} give
\[
 e<\frac{(p-1)p^{2n}}{p^n-1}\le pz=z_1<pe.
\]
It is therefore central, $p\mid z_1$, and $p\nmid(p+1)e-z_1$. Since $n-1\ge2$, Theorem~\ref{thm:central-p2-rigidity} applies to the paired quotients and gives equal traces. Their isogenies are defined over $K$, so the original traces agree as well, the final contradiction.
\end{proof}

\begin{proof}[A second proof of Theorem~\ref{thm:prime-index-upper}]
Suppose $n\ge3$, $p\nmid e$, and $e<p^n$. This numerical range is contained in \eqref{eq:ss-prime-range}, so Theorem~\ref{thm:ss-prime-range} treats the comparison of two whole supersingular groups. For the full height-two filtration class, pass to the completed maximal unramified base. A determinant has the form $\chi_p^d\epsilon$, with $d\in\{0,1,2\}$ and finite $\epsilon\in\mu_{p-1}$. Since $p\nmid e$, the wild cyclotomic order modulo $p^n$ is $p^{n-1}$. Equality of the determinant truncations therefore forces equal $d$ and equal finite factors.

For dimension zero the special character is the Brauer lift of the semisimplified residual character, and Cartier duality gives dimension two. In dimension one, the connected--\'etale filtration of a potentially ordinary piece descends by uniqueness, so the possibilities are a whole supersingular group or two height-one graded pieces. For two filtered groups, $\chi_p\equiv1\pmod p$ makes the special character recoverable from the residual semisimplification. These are the structural reductions detailed in the proof of Theorem~\ref{thm:second-wild-exact}. A mixed filtered/supersingular comparison would require $e>p^{2n-3}$ by Proposition~\ref{prop:mixed-obstruction}; this contradicts $e<p^n\le p^{2n-3}$. The supersingular theorem handles the remaining case, proving $\nuMin(p,e,2)\le n-1$.

For $p<e<p^3$, Theorem~\ref{thm:general-arc} gives unequal special traces with isomorphic $p^2$-torsion, while the upper bound applies to $p^3$-torsion. For $p^3<e<p^4$, it gives failure at level $p^3$, while the upper bound applies at level $p^4$. These are exactly the two asserted index values.
\end{proof}

\begin{proof}[Proof of Corollary~\ref{cor:prime-recovery-intro}]
Write $S_n=(p-1)p^{2n}/(p^n-1)$. For $n\ge4$,
\[
 S_n<p^{n+1}\le p^{2n-3}.
\]
The first inequality is equivalent to $p^n>p$, and the second to $n\ge4$. Under $e<S_n$, Theorem~\ref{thm:ss-prime-range} excludes an unequal-character supersingular comparison. Proposition~\ref{prop:mixed-obstruction} excludes a mixed comparison because it would require $e>p^{2n-3}$. The determinant, dimension and filtered reductions in the preceding proof impose no further upper bound on $e$. Thus the same argument proves recovery throughout the full filtration class. For $n=3$ the preceding theorem gives recovery under $e<p^3$.

Since $e$ is a multiple of $p-1$, the identity
\[
 \frac{S_n}{p-1}=p^n+1+\frac{1}{p^n-1}
\]
shows that $e<S_n$ is equivalent to $e\le(p-1)(p^n+1)$. The hypothesis $p\nmid e$ remains in force. The index-two interval follows from the preceding proof. The exact mixed construction of Corollary~\ref{cor:sharp-mixed-prime}, at level $p^3$, gives unequal traces whenever $e>p^3$; the new full-class bound at level $p^4$ supplies the matching upper bound up to $(p-1)(p^4+1)$. At level $p^4$, the mixed construction gives unequal traces whenever $e>p^5$; recovery at level $p^5$ supplies the matching upper bound up to $(p-1)(p^5+1)$. These are the asserted index-three and index-four intervals. The lower constructions allow the field of index $e$ to vary, whereas both upper bounds hold over every field in the specified base class.
\end{proof}

For example, at $p=5$ the index is three for $125<e\le2504$ and four for $3125<e\le12504$, subject in both cases to $4\mid e$ and $5\nmid e$. At level $p^3$, the mixed obstruction occurs inside the larger supersingular range, so the distinction between the two classes is necessary there.
